% Chapter 7: Forces in Everyday Life (complete sample chapter)
\chapter{Forces in Everyday Life}

\step{A Counterintuitive Opening}

\begin{mainline}
The base of a summer thundercloud is about two thousand meters above ground. If nothing impeded a raindrop falling from that height and gravity accelerated it all the way, it should reach your umbrella at
\[
v=\sqrt{2gh}\approx\sqrt{2\times9.8\times2000}\ \text{m/s}\approx 200\ \text{m/s}.
\]
What does two hundred meters per second mean? A handgun bullet leaves the muzzle at roughly two or three hundred. Rain falling at that speed would make every drop a little bullet and every rainy outing a trip into a shooting range.

Yet yesterday you walked in the rain. Drops pattered on your umbrella and felt cool on your face, nothing more. Meteorological observations give an ordinary raindrop's landing speed as only about seven meters per second, nearly thirty times less than free fall's predicted two hundred.

Where did that thirtyfold speed difference go? Who caught the drop halfway down? Neither cloud nor wind, but an invisible force that pushes on you every second: air resistance.

Now look indoors. You push a bookcase with all your strength and it does not move. Chapter 6 just said that zero acceleration requires zero net force, yet you are plainly pushing and your hands ache. Who supplies the force that pushes back? The floor appears to do nothing; it does not even move.

Raindrops, bookcases, the chair supporting you above the ground, the tires making you turn, and a taut musical string share one list of forces. Chapter 6 gave three general rules saying net force determines how motion changes. This chapter answers the postponed practical question: \emph{which forces do we encounter every day, and what is each one's character?}
\end{mainline}

\step{Make a Guess First}

Before seeing answers, place three bets. Wrong answers are more useful than right ones.

Question one. A bookcase rests on the floor. You push horizontally with fifty newtons and it stays still. What static friction does the floor exert? (a) Zero because it does not move; (b) less than fifty newtons; (c) exactly fifty; (d) more than fifty, otherwise it would slide backward.

Question two. A rubber band holding one coin stretches a centimeter. How far does it stretch holding three identical coins? (a) Still one centimeter; the band cannot recognize coins; (b) about two, since stretching gets harder; (c) about three, proportional to hanging weight; (d) nine, with an ever-increasing combined effect.

Question three. Large and small raindrops fall from one cloud. Which falls faster? (a) Equally fast, as Chapter 5 says heavy and light objects fall alike; (b) larger ones; (c) smaller ones; (d) whichever feels like it.

Question four. During hard braking, experienced drivers pump the brakes to keep the wheels from locking and skidding. What is the main benefit? (a) Less tire wear; (b) locked wheels give less braking force and lose steering; (c) cooler brake pads; (d) merely a driving habit without physical basis.

Write the three answers down. We will check them at the end, especially the first, whose answer first reveals to many people that a force can adapt to circumstances.

\step{Hands-on Experiment}

\begin{experiment}[A Rubber Band's Character and a Book on a Slope]
\textbf{Materials}: a rubber band of uniform thickness, a dozen or so identical coins, a ruler, a pen, paper, a thick book, and a flat board or sturdy cardboard sheet.

\textbf{Part one: how well does a rubber band follow rules?} Hang one end from a fixed point at a table edge, or hold it suspended with your left hand. Hang a coin from the lower end, measure the band's length, and record it. Add coins one at a time, recording length and calculating extension after each. Plot coin count horizontally and extension vertically. What appears? Points almost on a straight line through the origin, at least with few coins.

\textbf{Part two: overstretch it.} Add more coins, or pull by hand to two or three times its original length, then remove all coins. Measure its new unstretched length: it has \emph{grown longer} and will not fully shrink back. Add this portion to your graph. The line bends and does not return. The band's obedience has a limit.

\textbf{Part three: a book on a slope.} Place the thick book at one end of the board and slowly raise that end, steepening the slope. Record the approximate angle when the book \emph{just starts sliding}, using its spine to sketch a line on paper. Start it sliding again; once moving, it continues even when you lower the angle slightly. Just about to slide and already sliding are different states, with different frictional prices.

\textbf{Part four: crumpled and flat paper.} Take identical sheets, crumpling one tightly and leaving the other flat. Hold them at the same height and release together. The ball falls almost straight down while the flat sheet wavers and lags. Their weights and gravitational forces are identical; only frontal area differs. This is air resistance's cheapest home demonstration. In a class, also try replacing the flat sheet with one folded in half repeatedly and see how much faster it falls.

\textbf{Think}: why does the band not recover in part two? In part three, is friction greater just before sliding or while sliding? What would part four show on the airless Moon?
\end{experiment}

\step{The Old Model Runs into Trouble}

Chapter 6 left economical tools: three rules and net force determines momentum change. Carry them into a kitchen and you immediately hit three walls.

Wall one: \emph{the rules supply no inventory}. Net force is easy to write, but where do forces on a bookcase or raindrop come from? Gravity is one; what next? How strong is a chair's support or a rope's pull? Without an inventory, Newton's second law is an unloaded gun.

Wall two: \emph{some forces seem to read minds}. Push the bookcase gently and it rests; double your effort and it still rests. Old intuition says the agent issues a fixed amount of force. The floor does otherwise: you supply fifty newtons, it counters fifty; you supply eighty, it counters eighty. It asks only how much is \emph{needed} to prevent motion. A force pricing itself according to need directly contradicts the old idea of a fixed emitted force. This is a central intuition to repair.

Wall three: \emph{some forces get stronger as you move faster}. Cyclists know slow riding is easy, fast riding hard, and effort grows faster than speed. We dismissed Aristotle's claim that force sustains velocity, yet air resistance makes him seem half right: high speed indeed requires constant pedaling. The distinction is that you oppose neither motion itself nor its existence, but \emph{resistance growing with speed}. The old model has no place for that force. Add it to the list with a completely different price schedule.

These walls point to one need: not more general rules, but \emph{a profile for each force}, describing its source, direction, magnitude, and limits of cooperation.

\step{Inventing New Concepts}

Like inventorying a toolbox, introduce the common forces one at a time. Note each one's source, direction, and what determines magnitude.

\textbf{Gravity.} An old friend: Earth's attraction of an object. It points vertically down and near ground has magnitude approximately mass times \(g\approx9.8\) newtons per kilogram. It is the list's only indiscriminate member, always present near Earth. Chapter 8 reveals it as the same thing that governs the Moon's motion.

\textbf{Normal force.} The support from chair to person, table to cup, and floor to bookcase. It comes from \emph{tiny contact-surface deformations}. Sitting depresses the chair slightly, and the material pushes back like countless compressed microscopic springs. The direction is always perpendicular, or normal, to the surface; its magnitude is often written \(N\). That magnitude is set by what prevents you from passing through the surface, not by its own choice. Seat an elephant and it supports the elephant's weight until the chair breaks.

\textbf{Tension.} The force of ropes, cables, and musical strings. This too comes from deformation: the rope stretches imperceptibly, pulling molecules apart so they pull back. It always acts along the rope and only pulls, never pushes; a flexible rope cannot shove. Tension is equal everywhere along a light rope, a convenience of idealizing it as light.

\textbf{Elastic force.} The force of springs, rubber bands, and trampoline surfaces, deformation's most direct expression. More extension within limits gives a harder restoring pull. Part one of your experiment measured its character: force increases in step with extension, a straight line. We will soon write this direct behavior as a law.

\textbf{Friction.} Tangential dragging between contacting surfaces. Its profile needs two pages:

\emph{Static friction} acts before surfaces slide relative to each other. Its distinctive trait is \emph{supplying the amount needed}, starting from zero and going up to a limit roughly proportional to how tightly the surfaces press together, the normal force. It opposes the tendency to slide, not actual motion. What is it like? A negotiating partner with a budget ceiling: whatever you ask, it matches until funds run out and talks collapse. What is it \emph{unlike}? A spring, which responds to deformation rather than the needs of the situation; or gravity, which always turns up for work, whereas static friction does not appear unless needed.

\emph{Kinetic friction} acts once surfaces slide relative to one another. Its magnitude is roughly fixed and proportional to normal force; its direction opposes relative sliding. Part three hinted that its price is usually slightly below the static maximum, making an already sliding object harder to stop.

\textbf{Drag.} The resistance of a fluid, air or water, to an object passing through it. It opposes motion relative to the fluid and grows with speed: roughly proportional at low speed, to speed \emph{squared} at high speed, also depending on frontal area and fluid density. This is what catches raindrops.

Remember a counterintuitive example for later: \emph{friction can point along motion}. Walking feet push backward on the ground, whose static friction points \emph{forward}, propelling you. You walk because that negotiating partner is willing to pay. A box not slipping in an accelerating truck is likewise accelerated by forward static friction. Friction always opposes motion is a classic mistake: it opposes \emph{relative sliding or its tendency}, which need not oppose overall motion.

With the list in hand, these forces are at work everywhere. Tug-of-war rewards shoe-to-ground friction budgets rather than sheer strength; between teams of comparable mass, cleated shoes almost guarantee victory. Chopsticks lift a peanut through static friction at fingers-to-chopsticks and chopsticks-to-peanut. A trampoline returns elastic force saved in a large deformation. A cable-stayed bridge suspends hundreds of tonnes of roadway through its steel cables' tension. Learn their names and you will stop asking why a resting object does not move and ask instead who pushes it now and how the accounts balance.

Names must go onto paper. Chapter 6's free-body diagram now starts work: isolate the chosen object by setting a system boundary, draw nothing else, and represent only forces \emph{other objects exert on it} with arrows. Repeat one discipline: every listed force is an \emph{interaction} with its third-law partner. Floor pushes bookcase, bookcase presses floor; rope pulls you, you pull rope. Include only the half acting on your object; put the partner on the other object's diagram, never together to cancel. For the immovable bookcase:

\begin{center}
\begin{tikzpicture}[>=Stealth,thick]
  \draw[fill=gray!15] (0,0) rectangle (2.2,1.4);
  \node at (1.1,0.7) {Bookcase};
  \draw (-0.8,0) -- (3.0,0);
  \draw[->,red!70!black] (0,0.7) -- (1.0,0.7) node[midway,above]{$F_{\text{push}}$};
  \draw[->,blue!70!black] (2.2,0.35) -- (1.2,0.35) node[midway,above]{$f_{\text{static}}$};
  \draw[->,green!50!black] (1.1,1.4) -- (1.1,2.4) node[midway,right]{$N$};
  \draw[->,green!50!black] (1.1,0.7) -- (1.1,-0.6) node[midway,right]{$mg$};
\end{tikzpicture}
\end{center}

Vertically, normal force \(N\) and weight \(mg\) balance; horizontally, push \(F_{\text{push}}\) and static friction \(f_{\text{static}}\) balance. Net force zero, bookcase stationary, accounts settled. The diagram contains \emph{no} arrow for a forward motion tendency. Free-body diagrams show forces, not motion, and certainly not inertia. Adding one invented arrow is this book's most common drawing accident.

\step{The Model in One Sentence}

\begin{oneline}
Every force comes from a specific interaction and sets its own price: gravity always reports for work; support and static friction supply what is needed up to a limit; elasticity and tension respond to deformation; kinetic friction charges roughly a fixed rate; drag grows with speed. Learn their characters, find the net force, and leave the rest to Chapter 6.
\end{oneline}

\step{The Mathematical Translator}

Now write those characters as formulas, checking each with extreme cases immediately.

\textbf{Hooke's law: elastic force is proportional to deformation.}
\[
F=-kx .
\]
What it describes: the restoring force when a spring extends or compresses by \(x\). Why this form: your coin experiment gave a straight line through the origin, and this is its form. The spring constant \(k\) identifies stiffness; a stiff spring has large \(k\). The minus sign gives direction, not magnitude: pull right and force points left, always seeking the original length. When it holds: moderate extension. When it fails: experiment part two, when overextension bends the line and permanently lengthens the band; Hooke's law no longer recognizes it.

Check: \(x=0\) gives \(F=0\), as an unstretched, uncompressed spring should exert nothing. Double \(x\) and \(F\) doubles, matching your graph. Reverse direction by compressing rather than stretching: negative \(x\) gives positive \(F\), reversing the force while still restoring original length. At a tiny scale, a steel bar with large \(k\), compressed by only a hair's diameter, still gives substantial force. This is the microscopic identity of support.

Plotting part one's data gives this portrait of the law:

\begin{center}
\begin{tikzpicture}[>=Stealth,scale=0.9]
  \draw[->] (0,0) -- (5.6,0) node[below left]{Extension \(x\)};
  \draw[->] (0,0) -- (0,3.4) node[above left]{Pull \(F\)};
  \draw[thick,blue!70!black] (0,0) -- (3.6,2.7) node[right]{Slope \(=k\)};
  \draw[dashed,red!70!black] (3.6,2.7) .. controls (4.2,2.85) and (4.6,2.8) .. (5.2,2.4);
  \fill[red!70!black] (3.6,2.7) circle (1.5pt) node[below right]{Elastic limit};
  \node[red!70!black] at (4.9,3.1) {Overstretched: curve bends};
\end{tikzpicture}
\end{center}

The straight segment's slope is \(k\), unchanged for one spring and steeper for a stiffer one. Beyond the elastic limit, the graph bends without returning. The formula fails then, not merely becoming gradually inaccurate but \emph{entering a different game}.

\textbf{Static and kinetic friction: an inequality and an equality.}
\[
f_{\text{static}}\le \mu_{\text{static}} N,\qquad f_{\text{kinetic}}\approx \mu_{\text{kinetic}} N .
\]
These describe the \emph{upper limit} of tangential contact force and its approximate sliding price. The coefficient of friction \(\mu\) depends on the two materials and their condition, dry, oiled, or icy, and must be measured. Why an inequality for static friction? It supplies what is needed: the left side is actual force, the right the budget ceiling. The actual value comes from what prevents sliding, not directly from the formula. This chapter's easiest algebra mistake is treating \(f_{\text{static}}=\mu_{\text{static}}N\) as an equality, assuming the negotiator always spends its full budget.

Check: \(N=0\) makes both zero. Hold a book beside a blackboard without pressing and there is no friction. With no push on the bookcase, static friction is zero, not \(\mu_{\text{static}}N\); nothing justifies a spontaneous horizontal force. Reverse your push to the left and friction turns right, opposing sliding tendency without preferences. At the large extreme, a push greater than \(\mu_{\text{static}}N\) ends negotiations; sliding begins and the price changes to \(\mu_{\text{kinetic}}N\).

Look up coefficients to get numbers: rubber tires on dry asphalt have \(\mu_{\text{static}}\) around 0.7; wood on wood about 0.3; steel on ice below 0.05. Braking distances on dry ground and ice differ by tenfold. Winter accidents often reflect two bills from this inequality with different \(\mu\).

The inequality also supports a major automotive technology, anti-lock braking, ABS. With a \emph{rolling} wheel, the tire's ground-contact point is instantaneously at rest relative to the ground, so friction is static. Lock it into a skid and friction becomes smaller and kinetic, while steering fails. ABS releases and reapplies brakes more than ten times per second to keep wheels in the static-friction region just short of sliding, extracting near-maximum braking while retaining steering. It makes braking more \emph{knowledgeable} about friction, rather than merely harder.

\textbf{Drag: a price schedule that grows with speed.}
\[
F_{\text{drag}}\approx \tfrac12\,\rho\, C_d\, A\, v^2 \qquad(\text{at higher speeds}),
\]
Here \(\rho\) is fluid density, \(A\) frontal area, and \(C_d\) a shape coefficient, small for streamlined objects and large for plates. What it describes: to clear the air ahead at speed \(v\), how much air must you push aside each second, and how hard? Why \(v^2\)? The air mass encountered per second is proportional to \(v\), and the speed imparted to each parcel also to \(v\). Multiply them and the square appears. Check: \(v=0\) means no drag; a stationary umbrella has no wind account. Double speed and drag quadruples, so cycling at twenty versus forty demands fourfold effort, not double, matching your legs' memory. Double area and drag doubles: try opening an umbrella while riding, in a safe open space. The motor industry has burned tonnes of petrol over this formula. At 120 kilometers per hour, most of a family car's engine power fights air, so designers streamline bodies, reduce \(C_d\) from 0.4 to the low 0.2s, and test front and rear in wind tunnels. Every tenth saved in \(C_d\) translates into a real range figure.

\textbf{Terminal velocity: when the accounts balance.} A falling drop has downward weight \(mg\) and upward drag growing with speed. When drag equals gravity, net force and acceleration are zero and speed \emph{stops increasing}. The drop does not stop; it falls uniformly at that speed. Equate the forces:
\[
mg=\tfrac12\,\rho\, C_d\, A\, v_t^2
\quad\Longrightarrow\quad
v_t=\sqrt{\dfrac{2mg}{\rho C_d A}} .
\]
Check: greater mass means greater \(v_t\), so bigger drops fall faster, answering the third guess. Greater frontal area means smaller \(v_t\), the whole principle of parachutes. Without air, \(\rho=0\), the formula diverges: vacuum has no terminal speed, and feathers fall with hammers. The divergence honestly says the model no longer applies.

Estimate with real data: a raindrop of one-millimeter radius has mass about \(4\times10^{-6}\) kilograms, four milligrams, and frontal area \(A=\pi r^2\approx3\times10^{-6}\) square meters. Take \(C_d\) as \(0.5\) and air density \(\rho\approx1.2\) kilograms per cubic meter. Substitute:
\[
v_t=\sqrt{\dfrac{2\times4\times10^{-6}\times9.8}{1.2\times0.5\times3\times10^{-6}}}
\approx\sqrt{44}\approx 7\ \text{m/s}.
\]
Seven meters per second, precisely the measured value. One formula catches an entire rainfall from two thousand meters up.

Change the numbers for a skydiver: seventy kilograms and a face-down area about 0.7 square meters give \(v_t\) around fifty-five meters per second, nearly two hundred kilometers per hour. A free-falling skydiver cruises at highway speed. Opening the canopy raises area to tens of square meters and lowers \(v_t\) to around five meters per second, comparable to jumping from a two-meter platform. A parachute does not block gravity; it \emph{increases frontal area so air takes over the account sooner}.

\textbf{Circular motion: turning has a price.} Finally, apply the inventory to circles. Chapter 5 said velocity includes direction: even at fixed speed, changing direction changes velocity and therefore requires acceleration. An object traveling at speed \(v\) around a circle of radius \(r\) has instantaneous acceleration magnitude
\[
a=\frac{v^2}{r},
\]
always toward the center. Newton's second law immediately demands an inward \emph{net force} of \(mv^2/r\). This account is called centripetal force, but listen carefully: \emph{centripetal force is a job title, not a new member of the inventory}. What is the analogy like? A hospital's doctor on duty: the position can be filled by physicians from different specialties. What is it unlike? A new force species: nobody by that name appears on the roster. You cannot point to an object and add a centripetal force; you say which real force or forces are doing centripetal work now. Anyone can fill the post. For a turning car, tire-ground static friction does; for a washing machine's spin cycle, drum-wall normal force on water does. Drops escape holes because that support resigns, letting them fly tangentially, not because a centrifugal force throws them. For the orbiting Moon, gravity fills the role. Drawing centripetal force as another arrow alongside gravity and support pays twice for one job, the chapter's second major trap.

\begin{center}
\begin{tikzpicture}[>=Stealth,thick]
  \draw[gray!60] (0,0) circle (1.8);
  \fill (0,0) circle (1.5pt) node[below left=1pt]{Center};
  \fill[blue!70!black] (1.8,0) circle (2.5pt);
  \draw[->,blue!70!black] (1.8,0) -- (1.8,1.5) node[midway,right]{Velocity \(\bm v\) (tangential)};
  \draw[->,red!70!black] (1.8,0) -- (0.5,0) node[midway,above]{Acceleration \(\bm a\) (inward)};
  \fill[blue!70!black] (0,1.8) circle (2.5pt);
  \draw[->,blue!70!black] (0,1.8) -- (-1.3,1.8) node[midway,above]{\(\bm v\)};
  \draw[->,red!70!black] (0,1.8) -- (0,0.6) node[midway,right]{\(\bm a\)};
\end{tikzpicture}
\end{center}

Remember two things from the graph: velocity and acceleration do \emph{not} share a direction. One is tangential, the other inward, always perpendicular. Speed may remain unchanged while direction bends at every instant; bending costs \(mv^2/r\).

Check: \(v=0\) gives \(a=0\), no circling and no bill. Larger \(r\) gives smaller \(a\): at equal speed a sharp bend is more dangerous than a gentle one, just as driving intuition says. As \(r\to\infty\), \(a\to0\): an infinite-radius circle is a straight line, whose uniform motion has no acceleration. The model reduces smoothly to the straight case.

\begin{mathbridge}
\textbf{How a raindrop reaches constant speed.} Suppose low-speed drag is proportional to velocity: \(F_{\text{drag}}=bv\). Take downward as positive. Newton's second law becomes
\[
m\frac{\mathrm dv}{\mathrm dt}=mg-bv .
\]
The right side vanishes at \(v=mg/b\), the terminal speed \(v_t\). Set \(u=v-v_t\), giving \(m\,\mathrm du/\mathrm dt=-bu\), an exponential decay:
\[
v(t)=v_t\bigl(1-\mathrm e^{-t/\tau}\bigr),\qquad \tau=\frac{m}{b}.
\]
The characteristic time \(\tau\) sets the timescale: after about three to five \(\tau\), speed is close to \(v_t\). Check: at \(t=0\), \(v=0\); as \(t\to\infty\), \(v\to v_t\); as \(b\to0\), \(\tau\to\infty\), so no drag means never reaching a terminal speed. Again divergence expresses inapplicability. Quadratic drag uses the same reasoning with a nonlinear equation and a slightly different curve, but the conclusion persists: velocity approaches a finite \(v_t\) monotonically.

\textbf{Where centripetal acceleration comes from.} Translate two velocity arrows separated by a short \(\Delta t\) to one origin. The line between their tips points exactly inward as \(\Delta t\to0\). Use arc length \(v\,\Delta t=r\,\Delta\theta\) and the velocity triangle \(|\Delta\bm v|=v\,\Delta\theta\); divide and \(a=|\Delta\bm v|/\Delta t=v^2/r\) appears. This uses only velocity's vector nature, with no new assumption.
\end{mathbridge}

\begin{challenge}
\textbf{Why a friction coefficient is a makeshift number.} Under a microscope, two apparently smooth surfaces are mountains. Only a few peaks actually touch, giving a true contact area much smaller than the apparent area and roughly proportional to pressure. This happens to explain approximate \(f\propto N\), but the coincidence is fragile: contacts adhere, plow grooves, and vary with speed and temperature. Racing tires' effective \(\mu\) can exceed \(1\), and a thin liquid-water layer on ice greatly reduces \(\mu_{\text{kinetic}}\). Thus \(\mu\) resembles the slope of an empirical curve more than a fundamental material constant. A microscopic theory must account for adhesion, deformation, and fracture at every contact, still an active research field.

\textbf{Does centrifugal force exist?} From the ground, approximately inertial, a spinning drop leaves tangentially without any outward force, so centrifugal force does not exist. From a rotating frame, however, an object at rest in it appears to be pulled outward. To preserve Newton's equation form there, mathematicians name the tendency \emph{inertial centrifugal force}. It is a bookkeeping term arising from frame choice, not an interaction exerted by an object. Both accounts are consistent internally; mixing them causes errors. This is another reason to keep mathematical models and physical explanations distinct.
\end{challenge}

\step{The Model's Limits}

The inventory and calculations are done. Now attach an expiration condition to every formula.

\textbf{Hooke's law's boundary} was tangible in your experiment. Beyond the elastic limit a material deforms permanently, like a lengthened band or bent paper clip, or breaks. Microscopically, elasticity is atoms tending to return after a small separation, fundamentally an electromagnetic interaction, revisited with electromagnetism and materials. Rubber hardly even has a good linear region: it stretches several times its length with a strongly curved graph. Hooke's true home is metal springs with small deformations.

\textbf{The friction model's boundary} is broader. \(f=\mu N\) is a useful engineering approximation, not a natural law. The coefficient \(\mu\) varies with material pairing, contamination, humidity, temperature, and sliding speed. Static-to-kinetic transition is not a clean cut but a stick-slip vibration, heard in violin bows, chalk, and squealing brakes. Determine direction by relative sliding tendency, never simply object motion. Winter roads exhibit this model's limits: ice makes \(\mu\) plunge, prompting salt to lower freezing point and tire chains whose protrusions bite the surface, already partly beyond simple friction. The same sliding word names different accounts at different temperatures and surfaces.

\textbf{The drag model's boundaries} lie at both ends. Very slow, tiny objects such as dust and oil drops experience drag proportional to speed, not speed squared, in a viscosity-dominated regime. Near sound speed, strong air compression again demands another formula. The middle square-law range covers raindrops, skydiving, cycling, and cruising cars, enough to earn it a place in this chapter. Fluid mechanics can handle the two ends.

\textbf{The circular-motion account's boundary} is whether anyone will fill the job. On wet pavement, smaller \(\mu_{\text{static}}\) may leave the \(v^2/r\) demand beyond friction's budget, and the car slips outward. It is not flung; nobody pays enough, leaving a path combining tangential and inertial slipping motion. Engineers ingeniously raise the bend's outer edge, tilting the road's normal force and providing an inward horizontal component to share friction's bill. High-speed trains take bends with little slowing thanks to this superelevation.

\textbf{Support supplied as needed is limited} by structural strength: chairs have weight limits and bridges design loads. A constraint is not magic but compressed material repaying through elasticity. Fail to repay and it collapses.

\step{Explaining the Opening Phenomenon}

Now unpack the raindrop's account. It begins nearly at rest in the cloud and accelerates for the first second or two. Drag rises with speed squared and catches gravity within seconds. Thereafter net force is zero and the drop finishes the remaining two thousand meters at about seven meters per second. Acceleration occupies only the first dozen or so meters; ninety-nine percent is uniform motion. That is why rain from two thousand meters and two hundred feels almost identical on your face.

Was \(v=\sqrt{2gh}\) wrong? No. It is a model with a use-by condition: negligible drag. For a coin dropped from a table, drag is less than one percent and the formula shines. For a raindrop crossing two thousand meters, drag dominates and direct substitution chooses the wrong model. Chapter 1's warning that a model is not the world has just received a real soaking.

The same account explains everyday examples. A skydiver's terminal speed is about fifty-five meters per second before opening, then enlarged frontal area lowers \(v_t\) to around five, just enough for a standing landing. Cats can often walk away from falls of several floors while humans cannot because \(v_t\) contains mass divided by area, favoring smaller bodies. Cycling beyond thirty becomes difficult because drag grows quadratically and power requires multiplying by another speed, quickly reaching the body's limit. One price schedule governs falling from clouds to cats.

Check the three opening bets. Question one is (c): a stationary bookcase has zero horizontal net force, so static friction exactly matches fifty newtons, supplying neither too little nor too much. Choosing (a) confuses no motion with no force; choosing (d) retains the intuition that opposing force must be stronger. Question two is (c): one coin gives one centimeter, three give about three, Hooke's plain straight line, provided the band stays within its elastic limit. Question three is (b): in \(v_t=\sqrt{2mg/(\rho C_d A)}\), drop mass grows with radius cubed while area grows only with its square, so bigger drops fall faster. Drizzle drifts and storms sting because of this square root. Question four is (b): rolling contact is instantaneously stationary relative to ground and uses the larger static-friction budget; locking switches to smaller kinetic friction and loses steering control. ABS does exactly this management faster than any driver. Three correct answers mean your intuition has been upgraded; wrong guesses also pay, showing where the old model yields.

\step{Thought Experiments and Challenges}

\begin{thought}[Magnify an Ant a Hundredfold]
An ant falls from a skyscraper, brushes off dust, and continues; a horse falling from the same height meets a terrible end. Keep an ant's density and shape fixed but enlarge every dimension a hundredfold into a giant ant. What happens when it falls?

Calculate ratios: mass grows as size cubed, so a hundredfold enlargement means a millionfold weight. Frontal area grows only as the square, ten thousandfold. In \(v_t=\sqrt{2mg/(\rho C_d A)}\), numerator outgrows denominator, making \(v_t\) grow as the square root of size. The giant ant's terminal speed is tenfold greater. Landing is worse: momentum to be stopped grows cubically, while shell and leg strength, set by cross-sectional area, grows only quadratically. Fairy-tale giant insects would first be killed by physics, not heroes. Conversely, shrink an elephant to mouse size and a fall from a table leaves it unharmed.

The same account governs the sky: why does bird mass top out near a dozen kilograms? Wing lift grows roughly with area while weight grows with volume. Bigger birds find it harder to lift themselves, and the largest living flyers already approach this limit. Why must a high-altitude skydiver use a parachute while a swift can dive straight from a cliff into flight? One square root places life and death at its opposite ends. Every size change quietly rewrites force accounts, a principle recurring in Chapter 9 on energy and in material-strength chapters.
\end{thought}

\begin{puzzles}
\item A truck accelerates forward and a box in its bed does not slip, speeding up with it. Friction is its only horizontal force. Which way does static friction point? Hint: first ask what accelerates the box horizontally, then what would happen with a perfectly smooth truck bed. Tendency is the only key to direction.

\item A spring extends two centimeters under a five-hundred-gram weight. Cut it into equal halves and use one half to support the same weight. How much does it extend? Hint: imagine two half-springs in series, each extending one centimeter to give two in total. A shorter piece of the same spring is stiffer.

\item A car takes a fixed-radius bend in rain. The safe speed engineers specify is almost independent of its mass, assuming dry pavement and a known friction coefficient. Why does mass cancel? Hint: write the required account \(mv^2/r\) beside friction's budget \(\mu mg\) and see which letter appears in both.

\item During a washing machine's spin cycle, water leaves holes in the clothing and then travels nearly straight. Someone says centrifugal force throws it radially outward. Refute this with a free-body diagram. Hint: when force disappears, what remains is the object's instantaneous \emph{velocity}, tangential to the circle rather than radial.
\end{puzzles}

The inventory is complete, but its first member, gravity, has only a temporary near-ground rule, \(mg\). Why that rule? How is it related to the force making the Moon circle Earth? Who pays the Moon's centripetal bill? Next chapter follows this thread into the sky.
